\section{Spectral and spherical estimates for the support-profile argument}
\label{prof:sec:spectral}

The support-profile argument begins with estimates that involve only the
spectrum and unrestricted rotation orbits.  These estimates have two uses.
At positive observation times they provide fixed-rate overlap concentration
and exponential-moment bounds with arbitrarily prescribed exceptional-set
rates. At the
smallest times a more accurate comparison with the spherical model gives
covariance estimates and an independent lower-bound argument.

\subsection{Conventions and the observation law}
\label{prof:subsec:conventions}

The matrix $J$ has the GOE normalization: its independent off-diagonal
entries have variance $1/n$, and its diagonal entries have variance $2/n$.
The diagonal does not affect the spin law.  Write $\pi_{\mathrm c}$ for
uniform probability measure on $\{-1,1\}^n$, and $\pi_{\mathrm o}$ for
uniform probability measure on the sphere of radius $\sqrt n$.  Define
\[
 Z(h)=\int \exp\{x^TJx/2+h^Tx\}\,\dd\pi_{\mathrm c}(x),
 \qquad
 Z_{\mathrm o}(h)=\int \exp\{x^TJx/2+h^Tx\}\,\dd\pi_{\mathrm o}(x).
\]
Thus the partition functions use probability, rather than counting or
surface-area, normalization.  The spin law at field $h$ is $\mu_h$.
Empirical averages are denoted by $\av{a}=n^{-1}\sum_i a_i$; vector norms
are Euclidean, matrix norms without a subscript are operator norms, and
$\|\cdot\|_{\mathrm{HS}}$ is the Hilbert--Schmidt norm.

Given $J$, sample $X$ from $\mu_0$ and observe
\[
 h_s=sX+B_s,\qquad r=\sqrt s,\qquad k=nr^3,
\]
where $B$ is an independent standard Brownian motion in $\R^n$.
Completing the square in the Gaussian likelihood shows that the
conditional spin law is $\mu_s:=\mu_{h_s}$.  Its mean and covariance will
be denoted by $m_s^{\mathrm{post}}$ and $\Sigma_s$, respectively.

There are two probability laws for the disorder.  The ordinary law is
the centered GOE law.  The planted law is its size bias by
$Z(0)/\E Z(0)$, together with the Gibbs spin and observation.  Under the
present normalization $\E Z(0)=e^{n/4}$.  Under the
planted law, $X$ is uniform on the cube and
\[
 J=W+XX^T/n,
\]
where $W$ is an independent centered GOE matrix.  This follows directly
by completing squares in its Gaussian density.  A second exact identity
will be used repeatedly.  Fix the spectrum and write $J=U\Lambda U^T$,
randomizing the finite eigenframe choices so that $U$ belongs to the
proper rotation group.

In eigen-coordinates $c=U^Th_s$, the planted joint
law of $(U,c)$ has density
\[
 \frac{Z(Uc)}{Z_{\mathrm o}(Uc)}
\]
relative to Haar measure in $U$ and the spherical-model observation law
in $c$.  In this ratio both partition functions use the co-rotated
matrix.  Indeed the conditional density of $U$ is
$Z_U(0)/Z_{\mathrm o}(0)$, whereas the observation density contributes
$Z_U(Uc)/Z_U(0)$ times a rotation-invariant Gaussian factor.
Averaging $Z_U(Uc)$ over $U$ gives $Z_{\mathrm o}(Uc)$.  In particular,
the marginal law of $c$ is exactly the spherical observation law.

Constants may change from line to line.  All rate budgets in this
section are fixed before $n$ tends to infinity.  At a small scale, an
$o(1)$ error may also be made small by decreasing a fixed upper bound
on $r$; the time itself need not tend to zero.  Thus a statement with
exceptional mass $e^{-Bk}$, for every fixed $B$, means that the allowed
upper scale and the size threshold may depend on $B$ and on the desired
fixed accuracy.  Polynomial factors are absorbed by rate slack when
$k/\log n$ is sufficiently large.

\subsection{Resolvents at the upper edge}
\label{prof:subsec:spectrum}

Put
\[
 \ell=\log\log(n+1000),\qquad d_0=n^{-1/3}\ell,
 \qquad b_{\mathrm{sc}}(z)=\frac{z-\sqrt{z^2-4}}2.
\]
Let $\rho_{\mathrm{sc}}$ be the semicircle probability law on $[-2,2]$.
The next estimates are imposed on the \emph{ordinary} disorder.  Later
planted probabilities will be intersected with this spectral event;
the event need not have large planted probability.

\begin{lemma}[Spectral estimates]\label{prof:lem:spectrum}
There are events of ordinary probability tending to one on which,
uniformly for $d_0\le d\le D_*$ with any fixed $D_*<\infty$,
\begin{equation}\label{prof:eq1}
 \lambda_{\max}(J)\le 2+o(d_0^2),\qquad
 \frac1n\Tr(2+d^2-J)^{-1}
 =b_{\mathrm{sc}}(2+d^2)
   +o\!\left(\frac{d}{\sqrt{nd^3}}\right).
\end{equation}
For every fixed integer $p\ge2$, the trace of the $p$th resolvent
power is asymptotic to its semicircle value and is $O_p(nd^{3-2p})$.
There are order $nd^3$ eigenvalues at depth comparable to $d^2$ below
$2$, and at most $Cnd^3$ within an upper-edge interval of that width.
Moreover $\|J\|\le3$ and the empirical spectral measure converges to
$\rho_{\mathrm{sc}}$. Fixed-rank interlacings preserve the resolvent
estimates when the evaluation point stays at least a fixed positive
multiple of $d^2$ from every pole in both spectra.
\end{lemma}

\begin{proof}
The Gaussian--chi representation is the classical tridiagonal
construction; see \cite[Theorem~2.1 and Section~2.3]{DumitriuEdelman2002}
for the general beta-ensemble model.  In the present normalization,
successive orthogonal Gaussian elimination gives a tridiagonal matrix
with the same spectrum as $J$.  Its independent diagonal entries are
$N(0,2/n)$ and its adjacent entries are $\chi_{n-i}/\sqrt n$,
$1\le i<n$.  To see the independence, rotate the first column below
the diagonal onto the second coordinate.  Orthogonal invariance leaves
a fresh independent GOE matrix on the remaining coordinates, and the
procedure can then be iterated.

Let $A$ be the deterministic Jacobi matrix with zero diagonal and
adjacent entries $\sqrt{1-i/n}$, and write the tridiagonal matrix as
$A+\Delta$.  Centered entries of $\Delta$ are sub-Gaussian at scale
$n^{-1/2}$.  Gaussian norm concentration and the second moment of a
chi variable give adjacent-entry means
$O((n(n-i))^{-1/2})$.  For
\[
 R=(2+d^2-A)^{-1},\qquad a_i=(d^2+i/n)^{1/2},
\]
the positive walk expansion gives
\begin{equation}\label{prof:eq:jacobi-diagonals}
 (R^p)_{ii}\le C_p a_i^{1-2p},\qquad p\ge1.
\end{equation}
Indeed the positive series
\[
 (R^p)_{ii}
 =\sum_{j\ge0}\binom{p+j-1}{j}
       (2+d^2)^{-p-j}(A^j)_{ii}
\]
is a sum over returning walks. Their count divided by $2^j$ is
$O((j+1)^{-1/2})$. A walk staying above $i/2$ gains
$e^{-cij/n}$ from the adjacent weights, so its contribution is bounded by
\[
 C_p\sum_{j\ge0}(j+1)^{p-3/2}
          e^{-c(d^2+i/n)j}\le C_p a_i^{1-2p}.
\]
By reflection and the binomial tail bound, the relative count of walks
that visit below $i/2$ is at most $Ce^{-ci^2/(j+1)}$. Their sum is at
most $C_p d^{1-2p}e^{-cdi}$. Since $nd^3\ge1$,
$di\ge i/(nd^2)$; exponential decay in this last quantity bounds the
sum by $C_p d^{1-2p}(1+i/(nd^2))^{(1-2p)/2}$, which is again
$C_pa_i^{1-2p}$.

The perturbation is small in the resolvent metric already at
$d=n^{-1/3}\sqrt\ell$:
\begin{equation}\label{prof:eq:relative-noise}
 \|R^{1/2}\Delta R^{1/2}\|=o(1)
 \quad\hbox{with probability tending to one}.
\end{equation}
The block-energy estimate below follows the partial-sum and
summation-by-parts method used in
\cite[Proposition~7 and Lemmas~8--9]{LedouxRider2010}; see also
\cite[Lemma~5.6 and Section~6]{RamirezRiderVirag2011}.
The resolvent precision needed here is verified explicitly.
With $b_i=\sqrt{1-i/n}$ and $b_0=b_n=0$, the Jacobi energy is
\[
 u^{\mathsf T}(2+d^2-A)u
 =\sum_{i<n}b_i(u_{i+1}-u_i)^2+
   \sum_i(2+d^2-b_i-b_{i-1})u_i^2.
\]
It controls
\[
 \sum_i a_i^2 u_i^2+
 \sum_{i<n/2}(u_{i+1}-u_i)^2.
\]
The identity
$2\Delta_{i,i+1}u_iu_{i+1}
=\Delta_{i,i+1}(u_i^2+u_{i+1}^2)
-\Delta_{i,i+1}(u_{i+1}-u_i)^2$
separates each noisy adjacent term into a diagonal contribution and a
squared difference.

The maximum entry of $\Delta$ is $o(1)$, so the
difference terms and the terms on the bulk half are harmless.  Divide
the first half into consecutive blocks of length comparable to $a_i^{-1}$,
merging a short last block.  The values of $a_i$ are comparable within
each block.  Sub-Gaussian maximal bounds, applied separately to each
shifted entry sequence, show that centered diagonal partial sums from
the beginning of a block are at most $\epsilon a_i$, except with
probability $Ce^{-c\epsilon^2na_i^3}$ per block.

For a block $B$ of length comparable to $a_B^{-1}$, summation by parts
against a centered diagonal sequence $\zeta_i$ with these partial sums
gives
\[
 \left|\sum_{i\in B}\zeta_i u_i^2\right|
 \le C\epsilon\left(
 a_B^2\sum_{i\in B}u_i^2+
 \sum_{i,i+1\in B}(u_{i+1}-u_i)^2\right).
\]
To see the scale, average the endpoint value of $u^2$ over the block
and bound the variation of $u^2$ by Cauchy--Schwarz. Multiplication by
the partial-sum bound $\epsilon a_B$ then gives the two terms displayed.
The first-half means are $O(n^{-1})$, a relative
$O((na_B^2)^{-1})=o(1)$ contribution to the potential term.

A dyadic shell contains $O(na_B^3)$ blocks. At
$d=n^{-1/3}\sqrt\ell$, the smallest value of $na_B^3$ is
$\ell^{3/2}$. Choose $\epsilon\downarrow0$ with
$\epsilon^2\ell^{3/2}\gg\log\ell$. The block counts and the exponents
then grow geometrically on successive shells, so the exceptions are
summable. Summing the block inequalities proves
\eqref{prof:eq:relative-noise}. Positivity of $R^{-1}-\Delta$ at this
value of $d$ also places the upper edge below
$2+n^{-2/3}\ell=2+o(d_0^2)$.

The relative perturbation estimate has placed the spectrum below the
resolvent pole. The trace estimate requires more precision: its error
must be small on the unnormalized trace scale. We first estimate the
linear and quadratic perturbation terms, then identify the deterministic
Jacobi trace through its Hermite equation.

Expanding the resolvent and using \eqref{prof:eq:relative-noise} gives
\[
 \Tr(2+d^2-J)^{-1}
 =\Tr R+\Tr R^2\Delta
   +O\bigl(\Tr R\Delta R^2\Delta\bigr).
\]
The quadratic remainder is nonnegative, so its expectation controls it
by Markov's inequality. Its centered-noise expectation
is $O(d^{-2})$, by \eqref{prof:eq:jacobi-diagonals}, Cauchy--Schwarz
for neighboring entries, and
$n^{-1}\sum_i a_i^{-4}=O(d^{-2})$.  For the mean contribution, first
use $R\preceq Cd^{-2}I$.  On the bulk half the adjacent-entry means
are $O(n^{-1/2})$ and the diagonals of $R^2$ are bounded; on the other
half the means are $O(n^{-1})$ and the trace bound applies.

The linear
term has standard deviation $O(d^{-2})$ and mean $O(d^{-1}+1)$.
These estimates extend to each fixed-ratio interval in $d$. In an
eigenbasis of $A$, with $\rho_i=(2+d^2-\lambda_i(A))^{-1}$, the quadratic
term is
\[
 Q(d)=\sum_{i,j}\rho_i\rho_j^2|\widetilde\Delta_{ij}|^2.
\]
It is nonnegative and decreases with $d$, regardless of the signs in
$\Delta$. For the linear term $L(d)=\Tr R^2\Delta$,
\[
 d\,\partial_d L(d)=-4d^2\Tr R^3\Delta.
\]
The diagonal bounds give this logarithmic derivative the same
$O(d^{-2})$ fluctuation scale as $L(d)$ on a fixed-ratio interval.
Integrating it controls the supremum there. Markov's inequality for
$Q$ and the corresponding second-moment bound for $L$, at threshold
$(nd^3)^{1/4}d^{-2}$, have summable failures over the dyadic intervals
because $nd^3$ starts at $\ell^3$ and grows geometrically. Finally
\[
 \frac{(nd^3)^{1/4}d^{-2}}{\sqrt{n/d}}
 =(nd^3)^{-1/4}\longrightarrow0,
\]
which gives the stated uniform trace error.

It remains to identify the deterministic trace.  Reversing the Jacobi
indices in its principal-determinant recurrence gives monic Hermite
polynomials.  The characteristic polynomial $P_n$ satisfies
\[
 n^{-1}P_n''-zP_n'+nP_n=0.
\]
Therefore $b_A=P_n'/(nP_n)$ obeys
\[
 b_A^2-zb_A+1=-b_A'/n.
\]
The right side is $O((nd)^{-1})$ by
\eqref{prof:eq:jacobi-diagonals}.  Continuation from large $z$ selects
the smaller branch and gives
$b_A(2+d^2)=b_{\mathrm{sc}}(2+d^2)+O((nd^2)^{-1})$.
Together with the perturbation estimate this proves
\eqref{prof:eq1}.

For a fixed integer $p\ge2$, take the $(p-1)$st finite difference at
spacing $hd^2$, with $\Delta_a f(u)=f(u+a)-f(u)$. The scalar identity is
\[
 \frac{(-1)^{p-1}\Delta_{hd^2}^{p-1}(t+d^2)^{-1}}
 {(p-1)!(hd^2)^{p-1}}
 =\prod_{j=0}^{p-1}(t+(1+jh)d^2)^{-1}
\]
For an eigenvalue $v$ of $J$, take $t=2-v$. The upper-edge bound
already proved in \eqref{prof:eq1} gives
$t+d^2\ge d^2-o(d_0^2)\ge d^2/2$ for all sufficiently large $n$,
uniformly for $d\ge d_0$. Hence
\[
 (1+2(p-1)h)^{-(p-1)}(t+d^2)^{-p}
 \le \prod_{j=0}^{p-1}(t+(1+jh)d^2)^{-1}
 \le (t+d^2)^{-p}.
\]
The same comparison holds for semicircle depths $t\ge0$. For fixed
$h>0$, \eqref{prof:eq1} gives the semicircle limit for the finite
difference with a relative error tending to zero. Letting $h$ decrease
after $n$ tends to infinity now squeezes the $p$th power to its
semicircle value.

For the edge counts at a shrinking scale, put
$t_i=(2-\lambda_i(J))/d^2$ and consider the finite measure
\[
 \nu_{n,d}=\frac1{nd^3}\sum_i(1+t_i)^{-2}\delta_{t_i}.
\]
To identify its limit, use the coordinate $x=(1+t)^{-1}$. The upper-edge
bound places the pushed-forward measures on $[0,1+o(1)]$, and for every
fixed integer $j\ge0$ their moments are
\[
 \int (1+t)^{-j}\,d\nu_{n,d}(t)
 =\frac{d^{2j+1}}n\Tr(2+d^2-J)^{-j-2}.
\]
The trace-power estimates include the total mass, corresponding to
$j=0$. Polynomial approximation on a fixed compact interval therefore
determines the weak limit of the pushed-forward finite measures. For
$d\to0$, the semicircle change of variables identifies its pullback as
\[
 \pi^{-1}\sqrt t(1+t)^{-2}\mathbf1_{\{t>0\}}\,dt.
\]
The interval $[1,2]$ has positive limiting mass and boundary of mass
zero. Since its weight $(1+t)^{-2}$ is bounded above and below, it
contains order $nd^3$ eigenvalues. This conclusion is uniform over
shrinking scales: any violating sequence would have the same moment
limits and hence the same interval-mass limit. The second-power trace bound
gives at most $Cnd^3$ eigenvalues in any fixed multiple of the
$d^2$ upper-edge window. For $d$ in a fixed interval bounded away
from zero, ordinary semicircle convergence gives the same conclusions
with constants depending on that interval. Applying the same
argument to $-J$ gives the lower-edge bound.

For the interlacing assertion, the two eigenvalue counting functions
differ by at most a fixed integer. Write $z$ for the evaluation point
and let $(z-cd^2,z+cd^2)$ be free of both spectra. The function
$(z-v)^{-p}$, defined outside this gap, has an extension across the
gap of total variation $O_p(d^{-2p})$. Integration by parts against
the difference of the counting functions therefore bounds the raw
trace difference by $O_p(d^{-2p})$. This argument also applies when
$z$ is interior to the full spectral range. For $p\ge2$ the error
is relative $O((nd^3)^{-1})$; for $p=1$ it is
$o(\sqrt{n/d})$, the unnormalized error scale in \eqref{prof:eq1}.
\end{proof}

We will use these estimates under conditional matrix laws through a
deterministic interlacing implication. If $E$ has codimension one,
$e\in E$, and
\[
 M=(P_EJP_E)|_E-ee^{\mathsf T},
\]
then the ordered eigenvalues satisfy
\[
 \lambda_j(J)\ge\lambda_j(M)\ge\lambda_{j+2}(J)
 \qquad(1\le j\le n-2).
\]
In particular $\lambda_{\max}(M)\le\lambda_{\max}(J)$, cumulative
upper-edge counts differ by at most two, and band counts differ by at
most four. Thus the retained spectrum of the actual full matrix supplies
the edge spread and trace estimates for its compressed block. The
first-trace error $O(d^{-2})$ is also $o(\sqrt{n/d})$ because
$nd^3\to\infty$. When a conditional block has a Haar eigenframe, its
concentration estimates can be proved uniformly on these inherited
spectra before intersecting with the original full-matrix spectral
event; this does not require conditioning on that event's probability.

\subsection{Conditioning a Gaussian on the sphere}
\label{prof:subsec:spherical}

For a fixed spectrum and field, choose the unique $\lambda>\lambda_{\max}(J)$
such that
\begin{equation}\label{prof:eq:spherical-root}
 D=(\lambda-J)^{-1},\qquad m_{\mathrm o}=Dh,
 \qquad \Tr D+\|m_{\mathrm o}\|^2=n.
\end{equation}
The left side is strictly decreasing from infinity to zero, so the
root is well defined.  The spherical Gibbs law is precisely
$N(m_{\mathrm o},D)$ conditioned on $\|x\|^2=n$, in the sense of
disintegration by its length density.  Superscript $0$ will denote
the zero-field quantities.

Lemma~\ref{prof:lem:spectrum} implies
\begin{equation}\label{prof:eq:zero-shell}
 |\lambda^0-2|\le Cd_0^2,\qquad \|D^0\|\le Cnd_0.
\end{equation}
Indeed at $2+Cd_0^2$ the normalized trace deficit is $O(d_0)$.
At distance $c/(nd_0)$ above the last pole, that pole contributes
enough to close this deficit if $c$ is small; monotonicity gives the
second bound as well as the first.

The density estimates used in this conditioning are recorded with
their proofs, since polynomial conditioning losses must be tracked
at the lower edge.

\begin{lemma}[Density at the mean]\label{prof:quadratic-mean-density}
Let $g$ be a finite-dimensional standard Gaussian vector and let
$Q=g^TAg+b^Tg$ with $A\succeq0$ and $\sigma^2=\Var Q>0$.
The density of $Q$ at $\E Q$ is at least $c/\sigma$, for an absolute
$c>0$.  The same lower bound, with a smaller absolute constant,
holds at a displacement $o(\sigma)$ along any sequence of such forms.
\end{lemma}

\begin{proof}
After diagonalization and normalization it is enough to consider
\[
 \sum_i a_i(g_i^2-1)+2\sum_i b_i g_i,\qquad a_i\ge0,
 \qquad 2\sum_i a_i^2+4\sum_i b_i^2=1.
\]
Argue by subsequences. If an $a_i$ stays bounded below, write the sum
as $a_i(g_i^2-1)+2b_ig_i+W$. Cantelli's inequality gives a positive
lower bound for $\Pprob(W\le a_i/2)$, and Chebyshev's inequality removes
$W<-C$ after a fixed large $C$ is chosen. On the remaining event the
one-coordinate target lies in $[-a_i/2,C]$, a compact interval strictly
above the minimum $-a_i-b_i^2/a_i$ of the quadratic. Its density there
has a positive lower bound. This proves the claim along the subsequence.

If $\max_i a_i\to0$ but $\|b\|$ stays bounded below, write
$g=z\,b/\|b\|+g_\perp$. The coefficient of $z^2$ tends to zero, and
the random correction to its linear coefficient has variance
$O((\max_i a_i)^2)$. Thus that linear coefficient stays bounded away
from zero with high probability. The remaining term has bounded second
moment. Restrict it to a fixed compact interval and the cross coefficient
to a small interval; the change of variables in $z$ then has a bounded
Jacobian and a preimage in a fixed compact interval, where the Gaussian
density is bounded below. In the remaining case Fourier inversion gives
a local Gaussian limit.  The quadratic-product formula supplies the
integrable bound
\[
 \bigl(1+4(\max_i a_i)^2t^2\bigr)^{-c/(\max_i a_i)^2}.
\]
The same proof permits an $o(1)$ displacement of the target in
standard-deviation units.
\end{proof}

Suppose $\lambda-2\asymp d^2$, where $d_0\le d\lesssim1$.
The shell variable has standard deviation comparable to $\sqrt{n/d}$.
The lower bound follows from $\Tr D^2\asymp n/d$, and its upper bound
uses \eqref{prof:eq:spherical-root} to give
$\|m_{\mathrm o}\|^2=O(nd)$ and
$m_{\mathrm o}^TDm_{\mathrm o}=O(n/d)$.
The mean-density lower bound controls the denominator of shell
conditioning. For a projection test we also need an upper bound for the
remaining length density. The spectral spread survives a compression
of fixed codimension.

More explicitly, write $x=m_{\mathrm o}+D^{1/2}g$, let
$S=\|x\|^2$, and let $Z=V^{\mathsf T}g$ for a matrix $V$ with a fixed
number of orthonormal columns. Then
\[
 p_{Z\mid S=n}(z)=p_Z(z)\,
       \frac{f_{S\mid Z=z}(n)}{f_S(n)}.
\]
Let $P=I-VV^{\mathsf T}$. After fixing $Z$, the quadratic matrix
for the remaining white coordinates is the compression
$\left.PDP\right|_{\operatorname{ran}P}$. The conditional covariance of
$x$ is $D^{1/2}PD^{1/2}$, whose nonzero eigenvalues agree with those of
this compression. The surviving spectral block of the compression and
the quadratic-product formula therefore give
$f_{S\mid Z=z}(n)\le C\sqrt{d/n}$ uniformly in the translated mean,
hence uniformly in $z$. The mean-density lemma gives the matching
lower scale for $f_S(n)$. This ratio supplies Gaussian moment and tail
bounds for every fixed collection of white projections. In particular,
\begin{equation}\label{prof:eq2}
 \left\|
 \E_{\mathrm{shell}}
 \left[D^{-1/2}(x-m_{\mathrm o})(x-m_{\mathrm o})^TD^{-1/2}\right]
 \right\|\le1+o(1).
\end{equation}
The density ratio gives uniform moment control; the sharp coefficient
requires the limiting conditional law.

To obtain the coefficient $1$,
take any sequence of unit projection directions.  The standardized shell variable and the projection have
joint Gaussian subsequential limits, since
\[
 \frac{\|D\|}{\sqrt{\Tr D^2}}
 =O((nd^3)^{-1/2})=o(1).
\]
Their possible limiting correlation comes only from the linear part
of the shell variable.  Fourier inversion conditional on the
projection, with the integrable bound just obtained, gives local
convergence at shell value zero against bounded projection tests.

The uniform projection tails make the squared projections uniformly
integrable, so this local convergence also gives their second moments.  Conditioning a
joint Gaussian at its mean, or at an $o(1)$ displacement, can only
decrease the second moment, which proves \eqref{prof:eq2} uniformly
over directions.

When $\lambda\le2+Cd_0^2$, the same density-ratio argument costs at
most $\ell^C$.  Indeed $D\preceq D^0$,
$\|m_{\mathrm o}\|^2=O(nd_0)$, and the edge levels and remaining
dyadic shells give
\[
 \Tr D^2\le
 C\bigl[(nd_0^3)(nd_0)^2+n/d_0\bigr],
 \qquad m_{\mathrm o}^TDm_{\mathrm o}=O((nd_0)^2).
\]
The spread at depth comparable to $d_0^2$ supplies the required
density upper bound after fixed-dimensional compression. Thus the
ratio to the lower bound at the mean is bounded by a power of $\ell$.

We will also condition coarse tail events on exact spheres. The rule
used below is the following radial comparison. For a fixed number of
Gaussian vectors, let $I_E(n,\ldots,n)$ be the unnormalized joint
density of their squared lengths at $(n,\ldots,n)$, restricted to an
angular event $E$. Suppose the Gaussian precision and its product with
the mean have polynomial norms. If $E_{\rm buf}$ contains the radial
images of $E$ for $|s_j-n|\le n^{-C_0}$, then
\[
 I_E(n,\ldots,n)\le n^C\,\Pprob_{\rm G}(E_{\rm buf})
\]
for a fixed $C$ after $C_0$ is sufficiently large. Indeed, in polar
coordinates the Gaussian log density and the radial log Jacobian vary
by $O(1)$ on this interval, uniformly in the angular variables.
Integrating along those rays proves the inequality. Dividing by a
reference shell density bounded below by an inverse power of $n$
gives the corresponding normalized bound. At each use below, $E$ is
specified by quadratic or fixed-degree Gram tests on a polynomial
ball. Their polynomial derivative and inverse-tolerance bounds ensure
that a slight relaxation of the tests contains this radial buffer.

\begin{lemma}[Zero-field spherical normalization]\label{prof:lem:spherical-Z}
With ordinary probability tending to one,
\begin{equation}\label{prof:eq3}
 e^{n/4}n^{-C}\le Z_{\mathrm o}(0)\le e^{n/4+o(n)}.
\end{equation}
\end{lemma}

\begin{proof}
The Gaussian shell formula gives
\[
 \log Z_{\mathrm o}(0)
 =\frac{n(\lambda^0-1)}2
  -\frac12\log\det(\lambda^0-J)
  +\log\frac{\text{length density for }N(0,D^0)\text{ at }n}
                 {\text{length density for }N(0,I)\text{ at }n}.
\]
The last term is $O(\log n)$ in the direction needed for the lower
bound.  Replacing $\lambda^0$ by $2+Cd_0^2$ in the first two terms
increases them by only $O(nd_0^3)$, by convexity and
\eqref{prof:eq1}.  At this new argument, concavity of log determinant
gives the upper bound
\[
 \log\det(z-J)\le\log\det(z-A)-\Tr(z-A)^{-1}\Delta.
\]
The linear correction has mean $O(1)$ and standard deviation
$O(\sqrt{\log n})$, by \eqref{prof:eq:jacobi-diagonals}.

Integrating the deterministic Hermite transform estimate shows that
the deterministic log determinant differs from the semicircle value
by $O(\log n)$.  Beyond a fixed large $z$, the same Hermite equation
and $\|A\|\le2$ give normalized-transform error
$O(n^{-1}(z-2)^{-3})$, so the integral to infinity is harmless.
Finally,
\[
 \frac{z-1}{2}-\frac12\int\log(z-v)\,\dd\rho_{\mathrm{sc}}(v)
\]
has value $1/4$ at $z=2$, as follows by integrating
$b_{\mathrm{sc}}$, and is increasing thereafter.  This proves the
lower bound.  The upper bound follows from the same Gaussian
representation at a fixed positive shift, using the length-density
upper bound and then sending the shift to zero.
\end{proof}

\begin{lemma}[Spherical observations]\label{prof:lem:spherical-observation}
Assume the spectral event in Lemma~\ref{prof:lem:spectrum}.
For every fixed tolerance $\epsilon>0$, when $k\gg\log n$ and the
allowed upper bound on $r$ is sufficiently small, the spherical
observation satisfies
\[
 \lambda=2+(1\pm\epsilon)s,
 \qquad \|m_{\mathrm o}\|^2=(1\pm2\epsilon)nr,
\]
except with probability $e^{-c_\epsilon k}$. For every fixed $B<\infty$ there
are constants $c_B,C_B>0$ such that
\[
 c_Bs\le\lambda-2\le C_Bs,
 \qquad \|h_s\|\le C_Br\sqrt n,
\]
except with probability $e^{-Bk}$. Polynomial conditioning factors
can be accommodated already at $k>K\log n$, with $K$ chosen after the
tolerance and rate budget. For every deterministic actual time
$0\le t\le s=r^2$, the same upper test gives
\[
 \lambda(h_t)-2\le C_Bs,\qquad \|h_t\|\le C_Br\sqrt n
\]
outside probability $e^{-Bk}$, where $k=nr^3$ is the nominal scale.
The lower and sharp width bounds and the sharp mean estimate above
concern the actual time $s$.
\end{lemma}

\begin{proof}
Test \eqref{prof:eq:spherical-root} at $\lambda=2+\gamma s$.
Before conditioning, $X\sim N(0,D^0)$ and $h_s=rg+sX$ with an
independent white Gaussian $g$. Put $\delta=\lambda-\lambda^0$.
Since $k\gg\log n$, we have $r/d_0\to\infty$ and
$\delta=\gamma r^2(1+o_n(1))$ for fixed $\gamma>0$.
The partial fraction identity
\[
 D^2D^0=\delta^{-2}(D^0-D)-\delta^{-1}D^2
\]
together with \eqref{prof:eq1} and $\Tr D^0=n$ gives
\[
 \frac{\E\|D(rg+sX)\|^2}{nr}
 =\frac1{2\sqrt\gamma}+\frac1{2\gamma^{3/2}}
   +O_\gamma(r)+o_n(1),
 \qquad
 \frac{n-\Tr D}{nr}=\sqrt\gamma+O_\gamma(r)+o_n(1).
\]
The leading expressions agree at $\gamma=1$; their errors are small
in the stated order of upper scale and dimension.

The signal square is $s^2X^{\mathsf T}D^2X$. Subtract
$(s^2/\delta^2)(\|X\|^2-n)$, which has mean zero and vanishes on the
exact sphere. In a white eigen-coordinate with
$t=\lambda^0-\lambda_i(J)>0$, the remaining quadratic coefficient is
\[
 \frac{s^2}{t}\left(\frac1{(\delta+t)^2}-\frac1{\delta^2}\right)
 =-\frac{s^2(2\delta+t)}{\delta^2(\delta+t)^2}.
\]
The pole in $D^0$ has disappeared. The spectral trace bounds give
operator norm $O_\gamma(r^{-2})$ and squared Hilbert--Schmidt norm
$O_\gamma(n/r)$ for this quadratic matrix. The noise square has the
same bounds. Gaussian exponentiation at a deviation of order $nr$
therefore gives exponent
$\min\{(nr)^2/(n/r),(nr)/r^{-2}\}\asymp k$.
Conditional on $X$, the cross term is Gaussian with variance
$O(n/r)$ when the signal square is $O(nr)$. The radial conditioning
comparison preserves these rates after polynomial loss.
Testing on either side of $\gamma=1$ proves the sharp assertions.

For a prescribed rate, first choose a large fixed $L$. At a lower
trial point $2+\gamma r^2$, the $L^3k$ modes at depth comparable to
$L^2r^2$ contribute noise-square mean comparable to $nr/L$, whereas
the trace deficit is comparable to $n\sqrt\gamma r$. Choose
$\sqrt\gamma$ smaller than a fixed multiple of $L^{-1}$. A
fixed-fraction Gaussian small-ball bound on these modes then has
failure $e^{-cL^3k}$, uniformly in $X$ because translation decreases
the probability of a centered ball.

For the upper width, test at $2+L^2r^2$. The trace deficit is of order
$nLr$, while the noise and signal means are $O(nr/L)$ and
$O(nr/L^3)$. The noise quadratic matrix has norm
$O(L^{-4}r^{-2})$ and squared Hilbert--Schmidt norm $O(n/(L^5r))$;
the canceled signal matrix has the corresponding bounds
$O(L^{-6}r^{-2})$ and $O(n/(L^9r))$. Their Gaussian tail exponents at
the deficit scale are at least $cL^3k$. For an actual time $t\le s$,
write $h_t=\sqrt t\,g+tX$ and use this same trial point. Both noise
and signal prefactors decrease, so the upper-width estimate retains
its nominal $k$-scale rate. Choose $L$ from the requested rate, then
the lower trial constant $\gamma$, then the upper scale so that these
spectral windows are available. Finally choose the logarithmic
threshold to absorb the polynomial shell cost.
On the exact sphere $\|X\|=\sqrt n$; Gaussian norm
concentration applied to $h_t=\sqrt t\,g+tX$ gives the stated field
bound uniformly over the deterministic choices $0\le t\le s$.
\end{proof}

\section{Comparison on unrestricted rotation orbits}
\label{prof:sec:unconstrained}

In this section the matrix, field, and spherical mean rotate together.
The spectral data and the field in eigen-coordinates are fixed.
The comparison first controls unnormalized replica integrals.
A controlled Brownian argument then supplies the lower partition-function
tail needed to normalize those integrals.

\begin{proposition}[Unrestricted orbit comparison]
\label{prof:prop:unconstrained}
Assume Lemma~\ref{prof:lem:spectrum}, let $d=\sqrt{\lambda-2}$ be
comparable to $r>0$, and suppose $\|h\|\le Cr\sqrt n$ and
$k=nr^3\gg\log n$. Fix the width-comparison factors. For every fixed
$B$ and every requested fixed relative accuracy, the following holds
outside Haar mass $e^{-Bk}$ once the allowed upper bound on $r$ is
sufficiently small:
\begin{equation}\label{prof:eq4}
 \left|\log\frac{Z(h)}{Z_{\mathrm o}(h)}\right|\le o(k).
\end{equation}
For each fixed replica number $p$ and fixed symmetric zero-diagonal
$p\times p$ matrix $P$ of sufficiently small norm, chosen before the
upper scale and exceptional set, the same Haar bound holds for
\begin{equation}\label{prof:eq5}
 \log\E_{\mu_h^{\otimes p}}
 \exp\!\left\{
 \frac{s}{2}\sum_{a,b}P_{ab}(X^a-m_{\mathrm o})^T(X^b-m_{\mathrm o})
 \right\}
 \le C\Tr(P^2)k+o(k).
\end{equation}
The leading constant and the norm threshold depend only on the
width-comparison factors; the allowed upper scale may also depend on
the fixed replica data. The joint Haar--posterior probability that
a pair overlap differs from $\|m_{\mathrm o}\|^2/n$ by more than any
prescribed fixed tolerance times $r$ is at most $e^{-ck}$, with $c>0$
depending on that tolerance.

The same conclusions are available at $k>K\log n$ with explicit
polynomial losses. For any prescribed $\eta>0$, the errors $o(k)$ in
\eqref{prof:eq4} and \eqref{prof:eq5} may be replaced by
$\eta k+C_*\log n$. The unabsorbed exceptional masses have the form
$n^{C_*}e^{-B'k}$ for the partition and moment bounds and
$n^{C_*}e^{-c'k}$ for the joint pair bound. Here $C_*$ may depend on
the fixed replica data, accuracy, and preliminary rate budget.
Choose $B'>B$ and then $K$ to absorb these losses; for the pair bound
retain a smaller positive rate than $c'$. At fixed $K$ the logarithmic
loss is a chosen small fraction of $k$, rather than an error vanishing
with $n$. The exceptional-mass conclusions
also hold after weighting Haar measure by $Z(h)/Z_{\mathrm o}(h)$,
with an arbitrarily small decrease of the rate exponents.
\end{proposition}

\subsection{Replica integrals and their tails}
\label{prof:subsec:replica-integrals}

For independent uniform signs, write
$R_{ab}=(X^a)^TX^b/n$.  On polynomially small boxes near the identity
Gram matrix, their off-diagonal overlaps have probability at most
the corresponding spherical probability times
\begin{equation}\label{prof:eq6}
 n^{C_p}\exp\!\left\{C_pn\max_{a\ne b}|R_{ab}|^4\right\}.
\end{equation}
The spherical density is proportional to
$(\det R)^{(n-p-1)/2}$, with polynomial normalization obtained by
successive projections.  For signs, the sum Chernoff bound has an
analytic finite-dimensional Legendre expansion near zero, with
identity covariance.

Its quadratic and cubic terms agree with
$-\tfrac12\log\det R$: the only nonzero third products are triangles,
and each has expectation one.  The difference thus begins at fourth
order, which proves \eqref{prof:eq6}.  Conditional on the Gram matrix,
Haar rotation supplies exactly the same angular integral for both
priors.  The boxes can have any sufficiently small inverse-polynomial
side length.  Their Gram eigenvalues remain bounded below, and a
smooth Gram factor changes the energy and the additional tilts in
\eqref{prof:eq5} by a polynomially bounded amount times the box width.

Under the spherical Gibbs energies, pair overlaps are $O(r)$.
For $a\gg r$ up to a small fixed constant, the event
$\max_{i\ne j}|R_{ij}|\ge a$ costs at most a polynomial times
$e^{-cna^3}$, for every fixed number of replicas.  To prove this,
one of $(X^i+X^j)/\sqrt2$ and $(X^i-X^j)/\sqrt2$ has squared length
at most $n(1-a)$.  Translations decrease the relevant Gaussian
small-ball probability.  Add a negative square tilt of size
$\eta a^2$, with $\eta>0$ small.  By \eqref{prof:eq1}, the normalized
trace lost under this tilt is $O(r+\sqrt\eta a)$.  The logarithmic
Gaussian integral then gives the asserted $na^3$ exponent; exact
length conditioning costs only a polynomial.

For the tilted calculation write $\xi^a=X^a-m_{\mathrm o}$.
Before shell conditioning these are independent Gaussians of
precision $D^{-1}$. Choose the norm threshold on $P$ so that
$s\|P\|\le\theta(\lambda-\lambda_{\max}(J))$ for a fixed $\theta<1$.
This is possible uniformly under the width comparison. Factoring the
Gaussian integral leaves the determinant
$\det(I-sP\otimes D)^{-1/2}$ and a normalized centered Gaussian law
with covariance
\[
 \bigl(I_p\otimes D^{-1}-sP\otimes I\bigr)^{-1}.
\]
The precision remains positive. The determinant series has no linear
term because $\Tr P=0$, and the remaining terms cost at most
$Cs^2\Tr(P^2)\Tr(D^2)\le Ck\Tr(P^2)$.
The new covariance is bounded between fixed shifted resolvents.
Under this law $X^a=m_{\mathrm o}+\xi^a$; the preceding small-ball
argument for their sums and differences still applies because it is
uniform in translations. The joint length numerator is estimated
under this normalized correlated law, while the original single-shell
densities remain the denominators. The radial comparison gives the
stated polynomial conditioning cost. Without a tilt, Gaussian
quadratic and linear tails show concentration about
$\|m_{\mathrm o}\|^2/n$ within any fixed tolerance times $r$.

Consequently the annealed unnormalized products, divided by the
corresponding powers of $Z_{\mathrm o}(h)$ and restricted to all
overlaps in a sufficiently small fixed interval, have logarithmic
cost at most $Ck\Tr P^2+o(k)$.  On overlaps of order $r$, the cost
$nr^4=rk=o(k)$ in \eqref{prof:eq6} is negligible. For a fixed replica
number, choose the overlap cutoff $\eta$ so that
$C_pa^4\le(c/2)a^3$ for $0<a\le\eta$. On the rest of the small interval,
the negative $na^3$ exponent under the tilted Gaussian law then
absorbs the quartic prior-comparison error. The argument applies to repetitions in disjoint
replica groups. For the proposition's joint Haar--posterior pair
conclusion, take $P=0$. For any prescribed fixed tolerance, restrict
the pair to the complement of the interval centered at
$\|m_{\mathrm o}\|^2/n$ with that tolerance times $r$ as its radius.
The no-tilt concentration estimate above, combined with the same
prior comparison, retains a strict negative rate of order $k$ for
this restricted pair integral.

Only the pair calculation requires a completely unrestricted
unnormalized overlap tail.  For every fixed $\eta>0$,
$|R_{12}|>\eta$ contributes at most $e^{-c_\eta n}$ relative to
$Z_{\mathrm o}(h)^2$. The field costs $O(nr)$; we will absorb it by
choosing the upper scale after $\eta$.
After rotating the orthogonal sum and difference, their angular
energy is at most
\[
 \exp\!\left\{
 \frac n4\bigl[(1+R_{12})^2+(1-R_{12})^2\bigr]+o(n)
 \right\}.
\]
At a parallel endpoint only the nonzero direction is needed.
Here the elementary one-direction bound is that a uniform direction
of squared radius $n$, tilted by strength $u\ge0$, has logarithmic
integral per coordinate at most $u^2/4+o(1)$, uniformly on compact
$u$-ranges, also in a fixed-codimension compression.

Indeed use a
centered Gaussian with precision $z-uJ$, $z>2u$, to obtain the upper
bound
\[
 \frac{z-1}{2}-\frac12\int\log(z-uv)\,\dd\rho_{\mathrm{sc}}(v).
\]
At a fixed shift the shell-density loss is subexponential.  For
$u\le1$, choose $z=1+u^2$; for $u\ge1$, send $z\downarrow2u$.
The derivative of the latter continuation is
$1-1/(2u)\le u/2$, proving the claimed quadratic upper bound.
Iterating over the two directions proves the display. Its exponent is
$n/2+nR_{12}^2/2+o(n)$. For the denominator we need only
$\log Z_{\mathrm o}(0)\ge n/4-o(n)$, which follows deterministically
from the retained spectral event. Indeed \eqref{prof:eq:zero-shell}
gives $\lambda^0\to2$ and a polynomial bound on $\|D^0\|$.
The zero-field shell formula and the density-at-mean lower bound give
\[
 \log Z_{\mathrm o}(0)\ge
 \frac{n(\lambda^0-1)}2
 -\frac12\log\det(\lambda^0-J)-O(\log n).
\]
For every fixed $\delta>0$, eventually $\lambda^0\le2+\delta$,
so the log determinant is at most $\log\det((2+\delta)I-J)$.
Semicircle convergence identifies its normalized limit. Sending
$\delta\downarrow0$ gives the claimed lower bound. Thus no additional
normalization event is needed here.

By symmetry $Z_{\mathrm o}(h)\ge Z_{\mathrm o}(0)$, so the denominator
removes the $n/2$ term. The Cram\'er rate of the overlap of independent
signs exceeds $R_{12}^2/2$ by a fixed positive amount on
$|R_{12}|\ge\eta$. Choosing the upper bound on $r$ so that the field
cost is smaller than this gap proves the global pair estimate.

\subsection{The lower partition-function tail}
\label{prof:subsec:controlled-rotation}

Set $F(U)=\log(Z(h)/Z_{\mathrm o}(h))$ on the rotation orbit.
The Haar averaging behind this first-moment identity follows
Comets \cite[Equations~(2.1)--(2.2)]{Comets1996Spherical}; the present
field-dependent identity follows from the rotation convention above.
Haar averaging and the pair bounds give
\[
 \E_Ue^F=1,\qquad \E_Ue^{2F}\le e^{o(k)}.
\]
Paley--Zygmund therefore supplies a set of high anchors of mass
$e^{-o(k)}$ on which $F\ge-o(k)$.  These anchors may also be
required to satisfy
\[
 \left\|\E_{\mu_h}XX^T\right\|_{\mathrm{HS}}=o(n).
\]
Indeed the square of this Hilbert--Schmidt norm is
$\E_{\mu_h^{\otimes2}}(X^1\cdot X^2)^2$, and the global pair
estimate removes the contrary event at exponentially small
unnormalized mass.

For every fixed $B'$ and $\epsilon>0$,
\[
 \Pprob_U(F>\epsilon k)\le e^{-B'k}
\]
once the upper bound on $r$ is small enough. To see this, let
$q_\eta(U)=\mu_h^{\otimes2}(|R_{12}|>\eta)$. The global pair estimate is
$\E_U[e^{2F}q_\eta]\le e^{-c_\eta n}$. Hence
\[
 \Pprob_U(F>\epsilon k,\ q_\eta>\delta)
 \le\delta^{-1}e^{-2\epsilon k}\E_U[e^{2F}q_\eta].
\]
Fix a replica number $p$ and choose
$\delta<(2\binom p2)^{-1}$. On the remaining orientations at least
half of the posterior $p$-replica mass has all overlaps below $\eta$.
The annealed bound for that truncated integral then gives upper-event
mass at most $2e^{-(p\epsilon-o_p(1))k}$. Choose $p$ from $B'$, then
the truncation tolerance and the upper scale. This obtains the strong
upper tail from truncated moments.

We now turn the upper-tail bound into a lower-tail bound. The tangent
inequality at a high anchor compares nearby endpoints. A reflected third
rotation will cancel their first-order displacement, while its entropy
cost allows us to control its endpoint by the upper tail just proved.

At a high anchor, Jensen's inequality gives the deterministic lower
tangent estimate
\begin{equation}\label{prof:eq7}
 F(UR)\ge F(U)+\langle H_U,R-I\rangle_{\mathrm{HS}}
       -o(n)\|R-I\|_{\mathrm{HS}}^2,
 \qquad \|H_U\|_{\mathrm{HS}}=o(n).
\end{equation}
To verify it, expand the rotated quadratic and linear energies under
the Gibbs law at $U$.  The linear coefficient is bounded by a
constant times
$\|J\|\|\E XX^T\|_{\mathrm{HS}}+\sqrt n\|h\|=o(n)$.
The quadratic remainder is bounded by
$C\|J\|\|\E XX^T\|\|R-I\|_{\mathrm{HS}}^2$.
The denominator is rotation invariant.  This proves
\eqref{prof:eq7}, with arbitrarily small fixed leading errors at
sufficiently small upper scale.

Suppose, toward a contradiction, that the low set
$\{F<-\epsilon k\}$ has Haar mass at least $e^{-Bk}$.
The following controlled reflection produces three nearby rotations:
one at a high anchor, one in the low set, and one whose upper tail
can be bounded by entropy.

Use rotational Brownian motion started at the identity, with
right-translated antisymmetric noise $\dd W$ whose upper entries
have variance $\dd t/n$.  Its It\^o drift is
$-(1-1/n)I/2$.  Choose a deterministic terminal time sufficiently
large that its endpoint law approximates Haar as accurately as
needed.  The heat kernel on the compact connected group permits
this even for the exponentially small target sets at each fixed
$n$.

The control cost is the entropy--drift representation associated
with F\"ollmer's process; the Euclidean form and its converse are
given in \cite[Proposition~1 and Theorem~2]{Lehec2013Entropy}.
The group normalization needed here follows directly from the
likelihood martingale.  For a positive smooth terminal likelihood of free mean
one, its backward heat extension is the density-process function.
It\^o's formula under the Doob transform shows that the expected
integrated squared drift in standard noise coordinates is twice
the endpoint relative entropy.

An antisymmetric matrix velocity
$E$ corresponds to standard-coordinate drift with squared norm
$n\|E\|_{\mathrm{HS}}^2/2$.  Thus its entropy cost is exactly
\begin{equation}\label{prof:eq:control-entropy}
 \frac n4\E\int\|E_t\|_{\mathrm{HS}}^2\,\dd t.
\end{equation}
Bounded positive smooth approximations to normalized target
indicators attain endpoint hit probabilities arbitrarily close to
one and entropies within an additive constant of the indicator
entropy.  Conversely Girsanov's formula and data processing bound
the endpoint entropy of any adapted control above by
\eqref{prof:eq:control-entropy}; stopping at finite energy justifies
this statement without further integrability assumptions.

Choose feedback controls for two motions $U,V$ so that their
endpoints hit the high-anchor and low sets, respectively, with
probability at least $0.95$.  Their expected joint matrix energy is
$O_B(k/n)$.  Stop all controls when the joint energy exceeds a
sufficiently large fixed multiple of $k/n$.  Markov's inequality
makes this a small probability loss, so the two endpoint requirements
still hold jointly with probability greater than $0.7$.

Couple their noises as follows, and introduce a third motion $V'$.
If $D=U^TV$, conjugate the noise for $V$ to
$D^{-1/2}\dd W D^{1/2}$; use the analogous conjugation for $V'$.
These are Brownian marginal noises because conjugation is
orthogonal on antisymmetric matrices.  In base coordinates the
three controls are
\[
 E_U,\qquad E_V,\qquad 2E_U-E_V.
\]
Here $E_V$ is obtained by undoing the conjugation on the prescribed
feedback drift for $V$, which preserves its norm.  The first two
motions therefore retain their prescribed marginals up to the
energy stopping.  Initially use the local square roots near the
identity; the estimates below show that their chart is never left.

The identity
\[
 \dd W\,S\,\dd W=(S^T-\Tr(S)I)\,\dd t/n
\]
gives the relative equation
\begin{align*}
 \dd D={}&-\dd W D+\sqrt D\,\dd W\sqrt D\\
 &+\left[-(1-1/n)D+
             \frac{\Tr\sqrt D}{n}\sqrt D-\frac In\right]\dd t
   +(\sqrt D E_V\sqrt D-E_UD)\dd t.
\end{align*}
The same equation holds for $D'=U^TV'$ with the reflected control.
The noise has trace zero, so $q_D$ below has finite variation. For
$q_D=\|D-I\|_{\mathrm{HS}}^2
=2(n-\Tr D)$, expansion in rotation angles gives
\[
 -(1-1/n)\Tr D+\frac{(\Tr\sqrt D)^2}{n}-1
 =\left(\frac14-\frac1{2n}
       +O(\|D-I\|^2)\right)q_D.
\]
The control contribution to $q_D$ is bounded by
$C\sqrt{q_D}(\|E_U\|_{\mathrm{HS}}+\|E_V\|_{\mathrm{HS}})$.

Hence, while in the chart,
\[
 q_D'\le-cq_D+C\sqrt{q_D}\,e_t
       \le-\tfrac c2q_D+Ce_t^2,
 \qquad e_t=\|E_U\|_{\mathrm{HS}}+\|E_V\|_{\mathrm{HS}}.
\]
The reflected motion obeys the same inequality with a constant
multiple of $e_t$.  The energy cutoff therefore gives
\begin{equation}\label{prof:eq:relative-radius}
 \sup_t\bigl(\|D_t-I\|_{\mathrm{HS}}
             +\|D'_t-I\|_{\mathrm{HS}}\bigr)
 \le C_B\sqrt{k/n}.
\end{equation}
Taking $r$ sufficiently small in terms of $B$ prevents the first
exit and verifies the use of the local square roots.

The two relative rotations have opposite first-order displacements.
Their sum therefore begins with the quadratic error:
\begin{equation}\label{prof:eq:reflection-error}
 \E\|D_t+D'_t-2I\|_{\mathrm{HS}}^2\le C_B(k/n)^2,
\end{equation}
uniformly in the deterministic time $t$. This uniformity is needed
when the terminal Brownian time is taken large to approximate Haar
measure. We retain the damping term in the error equation to obtain it.  Put $A=D-I$, $A'=D'-I$, and
$H=A+A'$.  The linear noise in $A$ is $[A,\dd W]/2$ and its linear
drift is $-(1/2-1/n)A$.

The term $(\Tr A)I/(2n)$ is quadratic
because $\Tr A=-\|A\|_{\mathrm{HS}}^2/2$.  The constant control
terms cancel upon adding the two equations.  If
$a^2=C_Bk/n$, the resulting equation is
\[
 \dd H=\tfrac12[H,\dd W]-(1/2-1/n)H\,\dd t
       +\dd M_R+R_t\,\dd t,
\]
where the Taylor remainders satisfy
\[
 \frac{\dd\langle M_R\rangle_{\mathrm{HS}}}{\dd t}\le Ca^4,
 \qquad \|R_t\|_{\mathrm{HS}}\le Ca^2+Ca e_t.
\]
These follow from \eqref{prof:eq:relative-radius} and the convergent
matrix power series for the square root.

For every real matrix $H$
(the sum here need not be antisymmetric), the conditional noise
variance is
\[
 \frac{\E\|[H,\dd W]/2\|_{\mathrm{HS}}^2}{\dd t}
 =\frac{1-1/n}{2}\|H\|_{\mathrm{HS}}^2
 +\frac{\Tr(H^2)-(\Tr H)^2}{2n}
 \le\frac12\|H\|_{\mathrm{HS}}^2.
\]
It\^o's formula and Young's inequality, with small fixed constants
to absorb the quadratic-covariation cross term, yield
\[
 \frac{\dd}{\dd t}\E\|H\|_{\mathrm{HS}}^2
 \le-c\E\|H\|_{\mathrm{HS}}^2
      +Ca^4+Ca^2\E e_t^2.
\]
Integrating against the exponential damping and using the pathwise
bound $\int e_t^2\,\dd t\le Ca^2$ proves
\eqref{prof:eq:reflection-error}.

It remains to put the endpoint events and the matrix error on one
positive-probability event. Let $P_T$ be the free endpoint law and
$Q_T$ the reflected endpoint law. By \eqref{prof:eq:control-entropy},
$\operatorname{KL}(Q_T\Vert P_T)\le C_Bk$. For an event $E$ with
$0<P_T(E)<1$, binary relative entropy gives
\[
 Q_T(E)\le
 \frac{\operatorname{KL}(Q_T\Vert P_T)+\log2}
      {\log(1/P_T(E))}.
\]
If $P_T(E)=0$, finite relative entropy gives $Q_T(E)=0$.
Choose the upper-tail rate $B'$ large relative to $C_B$, and choose
the terminal Brownian time so that its Haar approximation error is
smaller than $e^{-B'k}$. The upper-tail estimate then gives
$P_T(F>\epsilon k/2)\le2e^{-B'k}$, making the reflected failure
probability as small as required. Markov's inequality in
\eqref{prof:eq:reflection-error} likewise gives its stated order
with any desired fixed high probability.

There is consequently
positive probability of the two endpoint hits, this bound for $V'$,
and both relative-matrix bounds, simultaneously.  On that event,
sum \eqref{prof:eq7} from the high anchor $U$ to $V$ and $V'$.
The linear error is at most
$o(n)\|D+D'-2I\|_{\mathrm{HS}}=o(k)$, and the two quadratic
errors are $o(k)$ by \eqref{prof:eq:relative-radius}.
Thus $F(V)+F(V')\ge-o(k)$.  This contradicts
$F(V)<-\epsilon k$ and $F(V')\le\epsilon k/2$ after fixing the
little-oh accuracies.  The lower tail, and hence
\eqref{prof:eq4}, follows.

\begin{proof}[Completion of Proposition~\ref{prof:prop:unconstrained}]
Normalize the replica integral bounds from
Section~\ref{prof:subsec:replica-integrals} using \eqref{prof:eq4}.
To handle large overlaps between replica groups after normalization,
put $q_\eta(U)=\mu_h^{\otimes2}(|R_{12}|>\eta)$. The global pair
bound gives
\[
 \E_U[e^{2F}q_\eta]\le e^{-c_\eta n}.
\]
On the retained lower-partition event $F\ge-\epsilon k$, Markov's
inequality shows that $q_\eta\le e^{-c_\eta n/2}$ outside Haar mass
$e^{-c_\eta n/2+2\epsilon k}$. A union bound applies to any fixed
number of replica pairs. Their fixed tilts cost at most $C_pnr^2$,
since $\|X^a\|=\sqrt n$ and $\|m_{\mathrm o}\|\le\sqrt n$.
Choosing the upper scale after $\eta$ therefore makes the tilted
exterior exponentially small in $n$ as well.

On the remaining replica configurations the truncated integral
estimates apply. Markov's inequality with repeated tilted replica
groups gives \eqref{prof:eq5}, choosing the denominator accuracy
arbitrarily small first. The pair concentration estimate transfers
in the same way, taking the denominator error below its strict
rate. The polynomial Gram-box and shell-density factors cost
$C_*\log n$ in these logarithmic estimates and in the exponents
of the raw failure bounds. Choosing $K$ after the rate
and accuracy absorbs them when $k>K\log n$, as stated.

Finally, for an event with an arbitrarily prescribed raw rate,
\[
 \E_U[e^F\mathbf1_E]\le
       (\E_Ue^{2F})^{1/2}\Pprob_U(E)^{1/2}
\]
gives any desired weighted rate by starting with more than twice that
rate. For a bounded posterior failure probability $0\le q(U)\le1$
having only a fixed strict rate, split instead at $F=\delta k$:
\[
 \E_U[e^Fq]\le e^{\delta k}\E_Uq+
       (\E_Ue^{2F})^{1/2}\Pprob_U(F>\delta k)^{1/2}.
\]
The second term has any fixed rate by the strong upper tail, and the
first loses only $\delta$ from its rate. Taking $\delta$ small proves
the stated weighted conclusions.
\end{proof}

\section{The polylogarithmic edge window}
\label{prof:sec:polylog}

The unrestricted comparison above has exponential accuracy on the
$k$-scale.  In the lowest spectral window a different second-moment
calculation gives convergence of the partition-function ratio to
one and an absolute error smaller than the critical covariance
scale.  This calculation also supplies an independent, supplementary
lower-bound argument.  At zero field the partition-function
comparison follows the quantitative critical comparison of
\cite[Theorem~1.6 and Section~3.3]{DuHuang2026CriticalFluctuations}.
The mixed cube--sphere cancellation for matrix moments is the
tensor-barycenter method of
\cite[Section~3.2, Lemma~3.5]{DuHuang2026CriticalOverlap}.
The field range and the weighted moment errors required here are
proved below.

\begin{proposition}[Edge comparison]\label{prof:prop:polylog}
Assume the spectral event of Lemma~\ref{prof:lem:spectrum}.  Suppose
\[
 r\le n^{-1/3}\log^{C_1}n,\qquad
 \lambda-\lambda^0\le n^{-2/3}\log^{C_2}n,\qquad
 \|h\|\le\sqrt n\,r\log n,
\]
where $C_1,C_2$ are fixed.  Set $L=n^{2/3}$ and $A_0=2I-J$.
Outside Haar mass $o(1)$,
\[
 \frac{Z(h)}{Z_{\mathrm o}(h)}=1+o(1).
\]
The cube and spherical raw second moments about $m_{\mathrm o}$
differ in operator norm by at most $n^{1/2+o(1)}$.  The cube raw
second moment, sandwiched by $A_0$, has operator norm at most
$n^{1/6+o(1)}$.  Each exponent bound can be read with any fixed
positive exponent slack.  Unspecified powers of logarithms may
depend on $C_1,C_2$.
\end{proposition}

\begin{proof}
For two independent prior vectors, the sphere--sphere and
cube--sphere overlap distributions agree exactly.  The cube--cube
overlap has lattice spacing $2/n$.  Stirling's formula and the
spherical density $c_n(1-R_{12}^2)^{(n-3)/2}$ show that its mass,
divided by the lattice spacing and the spherical density, equals
\begin{equation}\label{prof:eq:edge-Stirling}
 1+O(n^{-1/3}\log^C n)
 \qquad\text{for }|R_{12}|\le L\log^{C'}n/n.
\end{equation}
Choose $C'$ sufficiently large.  Outside this window the second
integrals, including any fixed bounded-degree polynomial insertions,
are negligible.

Indeed the negative-square-tilt argument of
Section~\ref{prof:subsec:replica-integrals} applies at depths much
larger than $(\lambda-2)_+^{1/2}+d_0$; combine it with
\eqref{prof:eq6}, then with the global pair estimate.  A cube-only
first moment averages exactly to the spherical first moment.

The following cancellation is needed to turn
\eqref{prof:eq:edge-Stirling} into an operator estimate.  Before
length conditioning, write independent Gaussian replicas as
\[
 x_i=m_{\mathrm o}+D^{1/2}\eta_i,\qquad i=1,2,
\]
with white $\eta_i$, and use Fourier inversion for
\begin{equation}\label{prof:eq:edge-statistics}
 \frac{\|x_1\|^2-n}{L},\qquad
 \frac{\|x_2\|^2-n}{L},\qquad
 \frac{x_1^Tx_2}{L}.
\end{equation}
The reciprocals of the two length densities at zero, in these
units, are at most $\log^C n$.

The Fourier inversions, including any fixed number of derivatives, cost
only powers of $\log n$ apart from the insertion size. Choose an edge
width $d_*=n^{-1/3}\log^M n$ larger than the polylogarithmic saddle
width. The spectral count supplies $N\asymp\log^{3M}n$ modes on which
$D_i/L\ge\log^{-C}n$. A frequency vector for the three statistics is
a symmetric $2\times2$ matrix $T$. At least one eigenvalue of $T$ has
size comparable to $\|T\|$, so the quadratic-product formula on those
modes gives
\[
 |\varphi(T)|\le
 \bigl(1+c\|T\|^2\log^{-C}n\bigr)^{-cN}.
\]
The noncentral factor does not increase the modulus, because the real
part of the inverse complex precision is positive. This bound is
integrable with any prescribed fixed frequency polynomial, with a
polylogarithmic integral. It remains so after deleting finitely many
modes, as required below.

We now assemble the signed comparison before estimating it. Let
$\pi_{\mathrm c}$ and $\pi_{\mathrm o}$ be the uniform cube and sphere
priors, put $E_U(x)=x^{\mathsf T}Jx/2+h^{\mathsf T}x$, and write
$H_U(x)=(x-m_{\mathrm o})(x-m_{\mathrm o})^{\mathsf T}-D$.
Set $Z_{\mathrm c}=Z$. For $i\in\{\mathrm c,\mathrm o\}$ define
\[
 z_i=\frac{Z_i}{Z_{\mathrm o}},\qquad
 \widehat M_i(U)=\frac1{Z_{\mathrm o}}
       \int H_U(x)e^{E_U(x)}\,d\pi_i(x),\qquad z_{\mathrm o}=1.
\]
All matrix data rotate with $U$. For a symmetric $A$ commuting with
$D$, expansion of the square gives
\[
 \mathbb E_U\|A(\widehat M_{\mathrm c}-\widehat M_{\mathrm o})A\|_{\rm HS}^2
 =\mathcal I_{\mathrm{cc}}-\mathcal I_{\mathrm{co}}
  -\mathcal I_{\mathrm{oc}}+\mathcal I_{\mathrm{oo}},
\]
where
\[
 \mathcal I_{ij}=\frac1{Z_{\mathrm o}^2}
 \iint\mathbb E_U\!\left[e^{E_U(x)+E_U(y)}
       \operatorname{Tr}\bigl(H_U(x)A^2H_U(y)A^2\bigr)\right]
       d\pi_i(x)d\pi_j(y).
\]
Conditional on the prior overlap, Haar integration gives the same
angular integral in all four terms. The three pair laws containing a
sphere use the same continuous overlap density; only the cube pair
uses the lattice and its multiplier in \eqref{prof:eq:edge-Stirling}.

After whitening, the common signed angular insertion is
\begin{equation}\label{prof:eq:edge-signed-kernel}
 \operatorname{Tr}\!\left[
 (\eta_1\eta_1^{\mathsf T}-I)K
 (\eta_2\eta_2^{\mathsf T}-I)K\right],
 \qquad K=D^{1/2}A^2D^{1/2}.
\end{equation}
Under the complex Fourier tilt its covariance differs from identity by
\[
 \Delta_T=(I-2\mathrm iT\otimes D/L)^{-1}-I.
\]
Its Hilbert--Schmidt norm, and the tilted mean using
$\|D^{1/2}m_{\mathrm o}\|/L\le\log^C n$, are bounded by frequency
polynomials times powers of $\log n$. Centering in
\eqref{prof:eq:edge-signed-kernel} cancels the baseline and the terms
belonging to only one replica. With zero tilted mean the remaining Wick
terms have the forms
\[
 \operatorname{Tr}(K\Delta_{22}K\Delta_{11}),\qquad
 \operatorname{Tr}(K\Delta_{21}K\Delta_{21}),\qquad
 (\operatorname{Tr}K\Delta_{12})^2.
\]
The first two are bounded using $\|K\|^2$ and the Hilbert--Schmidt
bounds on $\Delta_T$. For the last,
$|\operatorname{Tr}K\Delta_{12}|\le
 C\|T\|\operatorname{Tr}(|K|D)/L$. Terms containing the tilted mean
obey the same envelope. Fourier inversion therefore bounds the signed
density of the insertion and its overlap derivative by a power of
$\log n$ times
\begin{equation}\label{prof:eq:edge-kernel-size}
 \|K\|^2+\left(\frac{\operatorname{Tr}(|K|D)}L\right)^2.
\end{equation}
The frequency decay above absorbs all fixed derivative polynomials.

The spectral estimates give
\[
 \|D\|\le L\log^C n,\quad \operatorname{Tr}D^2\le L^2\log^C n,
 \quad \|A_0^2D\|\le\log^C n,\quad
 \operatorname{Tr}(A_0^2D^2)\le n\log^C n.
\]
The last two follow from $A_0D=I-(\lambda-2)D$ and $\|A_0\|\le5$.
Thus \eqref{prof:eq:edge-kernel-size} is at most $L^2\log^C n$ for
$A=I$ and $n^{2/3}\log^C n$ for $A=A_0$. Insert a smooth cutoff in
the central overlap window. The lattice spacing in the last statistic
of \eqref{prof:eq:edge-statistics} is $2/L$, so its Riemann-sum error
is the same envelope times $L^{-1}\log^C n$. The leading continuous
integrals cancel in the four-term sum, and
\eqref{prof:eq:edge-Stirling} leaves a squared error bounded by
$n^{-1/3}$ times the envelope, up to logarithms. The mass calculation
has envelope one.

Consequently the root-mean-square errors are
\[
 \begin{aligned}
 (\E_U|z_{\mathrm c}-1|^2)^{1/2}
   &\le n^{-1/6}\log^C n,\\
 (\E_U\|\widehat M_{\mathrm c}-\widehat M_{\mathrm o}\|_{\rm HS}^2)^{1/2}
   &\le n^{1/2}\log^C n,\\
 (\E_U\|A_0(\widehat M_{\mathrm c}-\widehat M_{\mathrm o})A_0\|_{\rm HS}^2)^{1/2}
   &\le n^{1/6}\log^C n.
 \end{aligned}
\]
Here $n^{-1/3}L^2=n$ and $n^{-1/3}n^{2/3}=n^{1/3}$ are the two
squared matrix scales. Markov's inequality gives the stated bounds
outside Haar mass $o(1)$, with any fixed exponent slack.

Finally let $Q_i$ be the normalized raw second moment about
$m_{\mathrm o}$. Then
\[
 Q_{\mathrm c}=D+\widehat M_{\mathrm c}/z_{\mathrm c},\qquad
 Q_{\mathrm o}=D+\widehat M_{\mathrm o},
\]
so the common $D$ cancels and
\[
 Q_{\mathrm c}-Q_{\mathrm o}
 =z_{\mathrm c}^{-1}(\widehat M_{\mathrm c}-\widehat M_{\mathrm o})
 +(z_{\mathrm c}^{-1}-1)\widehat M_{\mathrm o}.
\]
The shell density-ratio bound controls $\widehat M_{\mathrm o}$ in
operator norm by $L\log^C n$, and its $A_0$ sandwich by $\log^C n$.
The scalar mass error therefore costs at most $n^{1/2+o(1)}$ in the
unweighted comparison and less in the weighted one. The spherical
weighted raw moment itself is polylogarithmically bounded. This proves
all claims.
\end{proof}

\subsection{An independent covariance and slab lower bound}
\label{prof:subsec:supplementary-lower}

The realized-equilibrium-start and linear-test results in
\cite[Theorems~1.2--1.3]{AcceptedCritical} have stronger quantifiers:
they apply to every deterministic time sequence below the critical
scale and every deterministic diverging multiplicative slack,
respectively.  This supplementary subsection gives an independent
covariance obstruction and a worst-start slab proof with fixed
exponent slack.

At zero field, Proposition~\ref{prof:prop:polylog} and
\eqref{prof:eq3} give, with ordinary probability tending to one,
\begin{equation}\label{prof:eq:cube-Z-lower}
 Z(0)\ge n^{-C}e^{n/4}.
\end{equation}
This is also the lower partition-function event used later to
undo the planted size bias at polynomial cost.

On the same type of event the zero-field Gibbs covariance has an
eigenvalue at least $n^{2/3-o(1)}$.  Choose a spherical spectral
mode at depth of order $n^{-2/3}\ell^{10}$ below $2$, available by
Lemma~\ref{prof:lem:spectrum}.  Its unconditioned variance is
$n^{2/3}\ell^{-O(1)}$. More precisely, the chosen mode has variance
$O(n^{2/3}\ell^{-10})$, whereas the other modes at depth comparable
to $d_0^2$ give remaining length standard deviation at least
$c n^{2/3}\ell^{-1/2}$. Conditioning the chosen standardized
projection to lie between $1$ and $2$ therefore shifts the remaining
length target by $O(\ell^{-19/2})$ of that standard deviation.
Lemma~\ref{prof:quadratic-mean-density}
and the length-density upper bound show that this mode retains its
variance within an $\ell^{O(1)}$ factor on the sphere.

Symmetry at
zero field and the $n^{1/2+o(1)}$ cube--sphere error then give
\begin{equation}\label{prof:eq:supplementary-covariance}
 \|\Cov_{\mu_0}X\|\ge n^{2/3-o(1)}.
\end{equation}
For a unit linear test $f_v(x)=v^Tx$, the unscaled heat-bath form
satisfies
\[
 \D(f_v)=\sum_i\E_{\mu_0}
             \Var(f_v\mid X_{[n]\setminus\{i\}})
 \le\sum_i v_i^2=1.
\]
Thus the inverse unscaled spectral gap is at least
$n^{2/3-o(1)}$.  A top nonconstant eigenfunction gives the same
exponent lower bound for fixed mixing accuracies less than $1/2$.

The comparison also proves that worst-start distance tends to one
below this exponent.  Let $v$ be the same unit spectral mode, and
put
\[
 Y=\frac{v^TX}{\sigma_Y},\qquad
 \sigma_Y=n^{1/3}\ell^{-O(1)},
\]
where $\sigma_Y$ is its unconditioned spherical Gaussian standard
deviation.  For every fixed $K>0$ and $w=\ell^{-K}$, the spherical
probabilities of both $|Y|\le w$ and $|Y|\le2w$ lie between
$cw\ell^{-C}$ and $Cw\ell^C$.  The density estimates just used are
uniform when this standardized projection lies in the indicated
range, which proves these bounds.

For these events, the unnormalized cube probabilities, measured
in units of $Z_{\mathrm o}(0)$, differ from their spherical
counterparts by at most $n^{-c}$, outside ordinary Haar mass
$o(1)$, for an absolute small $c>0$.  This follows from the same
mass second-moment comparison, with the projection indicator on
each copy.  The angular integrals at fixed Gram matrix continue
to agree because $v$ co-rotates with the matrix.  In the central
overlap window, the bounded Stirling error is sufficient.

For
the lattice quadrature, condition first on the two tested white
Gaussian coordinates, each of absolute value at most one.
Deleting this one spectral mode leaves the eigenvalue spread for
the remaining two lengths and overlap.  Their joint density and
its overlap derivative are therefore bounded by powers of
$\log n$ in the units of \eqref{prof:eq:edge-statistics}, uniformly
in the conditioned coordinates.  The quadrature error is again
at most a power of $\log n$ divided by $L$.  The exterior of the
window has the same absolute tail bound as before.  Markov's
inequality proves the asserted $n^{-c}$ comparison.

Since the total mass ratio tends to one, choose $K>C+1$.
With ordinary probability tending to one, the narrow slab
$A=\{|Y|\le w\}$ then has Gibbs mass at least $n^{-o(1)}$, while
the wider slab $A^+=\{|Y|\le2w\}$ has mass tending to zero.
For the stationary rate-one-per-site chain, the spectral theorem
and $\D(f_v)\le1$ give
\[
 \E_{\mu_0}(f_v(X_t)-f_v(X_0))^2\le2t.
\]
For $m$ single-site attempts the same bound is $2m/n$, since
the attempt kernel is an average of conditional-expectation
projections and hence has spectrum in $[0,1]$.

Conditioning the stationary chain at its initial point to lie
in $A$ increases the bound by at most $1/\mu_0(A)$.
Consequently, for every fixed $\delta>0$ and
$t\le n^{2/3-\delta}$ in rate-one-per-site time,
\[
 \Pprob_{\mu_0(\,\cdot\mid A)}(X_t\notin A^+)
 \le\frac{2t}{\mu_0(A)\,w^2\sigma_Y^2}=o(1).
\]
The stationary mass of $A^+$ is $o(1)$, so this conditional-start
law has total variation distance $1-o(1)$ from equilibrium.
Averaging over its starting points, after the disorder and time have
been fixed, yields a starting configuration with the same conclusion.
In attempts, the corresponding time
range is $m\le n^{5/3-\delta}$.
